\textsc{[Greedy]}
\\
Implement Dijkstra's shortest path algorithm in C\texttt{++}.
Use adjacency list to describe the graph.
Assume the vertices are named 0, 1, 2, ..., $n - 1$
where $n$ is the number of nodes.
You also need to keep track on the \lq\lq distance from starting node to
vertex $i$", $d(i)$. You can use an array (you can use
$-1$ for infinity or a number larger than then longest possible path).
You also need to keep track on the weight from $i$ to $j$, $w(i,j)$.
You can either use a 2D array or a hashtable.
Ether one will achieve the runtime I mentioned in class (and in the notes),
although if the graph is huge and sparse a hashtable is probably better.
And for computing the solution, for each node $n$, you
also need to keep track of
\lq\lq current best previous node of $i$ to go back to the starting node",
$\operatorname{prev}(i)$.
You can use an array for $\operatorname{prev}$.
Note that you will need to use a modified minheap, i.e.,
the minheap comes with a hashtable that keeps track on the values in the
minheap.

Here's a sample execution:
\begin{Verbatim}[frame=single,commandchars=\\\{\}]
\userinput{4 5}
\userinput{0 1 10}
\userinput{0 2 11}
\userinput{1 2 200}
\userinput{1 3 100}
\userinput{2 3 12}
\userinput{0 3}
0 2 3 23
\end{Verbatim}
The first line of input \verb!4 5! is the number of nodes and number of edges.
In this case there are 4 nodes. They are named 0, 1, 2, 3.
And there are 5 edges.
The next 5 lines of output describes the weights in this graph.
For instance
\verb!0 1 10! says that the edge from node 0 to node 1 has weight 10,
\verb!0 2 11! says that the edge from node 0 to node 2 has weight 11, etc.
The graph is simple and undirected.
So the weight going from node 0 to node 1 is the same as the weight going from
node 1 to node 0.
The last line of input \verb!0 3! specifies the starting and ending node
for the shortest path.
The output
\verb!0 2 3 23! says that a shortest path is from 0 to 2 to 3
with path cost (i.e. sum of weights) of  $11 + 12 = 23$.
